\documentclass[10pt, aspectratio=169]{beamer}
% Preámbulo modular para presentaciones Beamer (Modelación Estadística)
% Diseño y paleta institucional del Tecnológico de Monterrey

\documentclass[aspectratio=169,xcolor={dvipsnames,table}]{beamer}

\usepackage[utf8]{inputenc}
\usepackage[spanish,mexico]{babel}
\usepackage{amsmath,amssymb,amsfonts,mathtools}
\usepackage{graphicx}
\usepackage{listings}
\usepackage{booktabs}
\usepackage{tikz}
\usetikzlibrary{matrix,arrows,decorations.pathmorphing,shapes.geometric,positioning}

% Paleta Institucional Tecnológico de Monterrey
\definecolor{TecAzul}{HTML}{0039A6}       % Azul TEC: color primario institucional
\definecolor{TecAzulOscuro}{HTML}{1A2E51} % Azul marino oscuro
\definecolor{TecRojo}{HTML}{EC2661}       % Rojo de acento (paleta miTec)
\definecolor{TecGris}{HTML}{646464}       % Gris neutro para comentarios y subtítulos
\definecolor{TecFondoBloque}{HTML}{F4F6F9}% Fondo suave para bloques

% Configuración de Tema Beamer
\usetheme{Boadilla}
\usecolortheme[named=TecAzulOscuro]{structure}
\setbeamercolor{alerted text}{fg=TecRojo}
\setbeamercolor{palette primary}{bg=TecAzulOscuro,fg=white}
\setbeamercolor{palette secondary}{bg=TecAzul,fg=white}
\setbeamercolor{palette tertiary}{bg=TecAzulOscuro!80!black,fg=white}
\setbeamercolor{title}{fg=TecAzulOscuro,bg=TecFondoBloque}
\setbeamercolor{frametitle}{fg=TecAzulOscuro,bg=TecFondoBloque}

% Bloques personalizados elegantes
\setbeamercolor{block title}{bg=TecAzulOscuro,fg=white}
\setbeamercolor{block body}{bg=TecFondoBloque,fg=black}
\setbeamercolor{block title alerted}{bg=TecRojo,fg=white}
\setbeamercolor{block body alerted}{bg=TecRojo!8,fg=black}
\setbeamercolor{block title example}{bg=TecAzul,fg=white}
\setbeamercolor{block body example}{bg=TecAzul!8,fg=black}

% Redondear bordes y añadir sombra sutil
\setbeamertemplate{blocks}[rounded][shadow=true]
\setbeamertemplate{navigation symbols}{}

% Pie de página limpio con número de diapositiva
\setbeamertemplate{footline}{
  \leavevmode%
  \hbox{%
  \begin{beamercolorbox}[wd=.7\paperwidth,ht=2.25ex,dp=1ex,left]{author in head/foot}%
    \hspace*{2ex}\insertshortauthor~~(\insertshortinstitute) \hfill \insertshorttitle
  \end{beamercolorbox}%
  \begin{beamercolorbox}[wd=.3\paperwidth,ht=2.25ex,dp=1ex,right]{date in head/foot}%
    \insertshortdate{}\hspace*{2em}
    \insertframenumber{} / \inserttotalframenumber\hspace*{2ex}
  \end{beamercolorbox}}%
  \vskip0pt%
}

% Configuración de Listings de Python para visualización pedagógica de código
\lstset{
    language=Python,
    basicstyle=\ttfamily\footnotesize,
    keywordstyle=\color{TecAzulOscuro}\bfseries,
    stringstyle=\color{TecRojo},
    commentstyle=\color{TecGris}\itshape,
    showstringspaces=false,
    breaklines=true,
    frame=lines,
    rulecolor=\color{TecAzulOscuro!30},
    backgroundcolor=\color{TecFondoBloque},
    aboveskip=0.5em,
    belowskip=0.5em,
    literate={á}{{\'a}}1 {é}{{\'e}}1 {í}{{\'i}}1 {ó}{{\'o}}1 {ú}{{\'u}}1
             {Á}{{\'A}}1 {É}{{\'E}}1 {Í}{{\'I}}1 {Ó}{{\'O}}1 {Ú}{{\'U}}1
             {ñ}{{\~n}}1 {Ñ}{{\~N}}1 {¿}{{?`}}1 {¡}{{!`}}1
}

% Metadatos institucionales por defecto
\title[Modelación Estadística]{Modelación Estadística para Ciencia de Datos e Ingeniería}
\author[J. Castillo Colmenares]{Juliho David Castillo Colmenares}
\institute[Tec de Monterrey]{Tecnológico de Monterrey \\ Escuela de Ingeniería y Ciencias}
\date{\today}

% Comandos y macros estadísticas compartidas para las presentaciones Beamer (Metropolis)

% Conjuntos numéricos básicos
\newcommand{\R}{\mathbb{R}}
\newcommand{\N}{\mathbb{N}}
\newcommand{\Q}{\mathbb{Q}}
\newcommand{\Z}{\mathbb{Z}}
\newcommand{\set}[1]{\left\{ #1 \right\}}
\newcommand{\abs}[1]{\left|#1\right|}
\newcommand{\norm}[1]{\left\| #1 \right\|}

% Comandos estadísticos y probabilísticos del curso
\newcommand{\Var}{\operatorname{Var}}
\newcommand{\cov}{\operatorname{Cov}}
\newcommand{\corr}{\rho}
\newcommand{\comb}[2]{\begin{pmatrix} #1 \\ #2 \end{pmatrix}}
\newcommand{\s}{\sigma}
\newcommand{\particion}{\mathfrak{P}}
\newcommand{\card}[1]{\#\left( #1 \right)}
\newcommand{\seq}[3]{\set{#1}_{#2}^{#3}}
\newcommand{\E}{\mathbb{E}}
\newcommand{\Prob}{\mathbb{P}}

% Entornos pedagógicos y cajas en español académico natural
\newenvironment{discusion}{%
  \begin{alertblock}{Pregunta de Discusión}%
}{%
  \end{alertblock}%
}

\newenvironment{reflexion}{%
  \begin{alertblock}{Reflexión para el Aula}%
}{%
  \end{alertblock}%
}

\newenvironment{paradoja}{%
  \begin{alertblock}{Paradoja de Exploración}%
}{%
  \end{alertblock}%
}

\newenvironment{motivacion}{%
  \begin{exampleblock}{Motivación: Ciencia de Datos en la Práctica}%
}{%
  \end{exampleblock}%
}

\newenvironment{retoclase}{%
  \begin{block}{Reto Rápido en Equipo}%
}{%
  \end{block}%
}


\title[Distribución Exponencial]{Sección 04.05: Distribución Exponencial}
\subtitle{Procesos sin Memoria, MLE, Distribución Erlang y Sistemas de Colas}
\author[J. Castillo Colmenares]{Juliho Castillo Colmenares}
\institute{Tecnológico de Monterrey}
\date{\vspace{-1.2cm}}

\begin{document}

% Slide 1: Portada
\begin{frame}[plain]
  \titlepage
\end{frame}

% Slide 2: Hoja de Ruta
\begin{frame}{Hoja de Ruta --- Capítulo 04: Variables Aleatorias Continuas}
  \footnotesize
  ¿Dónde nos encontramos en el desarrollo modular del capítulo? \pause
  \begin{itemize}\itemsep=0.08em
    \item \textbf{04.01} Función de Densidad (PDF) y Soporte Continuo \pause
    \item \textbf{04.02} Función de Distribución Acumulada (CDF) Continua y Cuantiles \pause
    \item \textbf{04.03} Esperanza, Varianza y LOTUS Continuo \pause
    \item \textbf{04.04} Distribución Uniforme Continua $U(a, b)$ \pause
    \item \textbf{04.05} Distribución Exponencial y Procesos sin Memoria \emph{(Hoy)} \pause
    \item \textbf{04.06} Distribución Normal y Puntaje $Z$ \pause
    \item \textbf{04.07} Distribuciones Gamma, Beta y Weibull
  \end{itemize}
  \vspace{-0.05cm}
  \begin{block}{Objetivo de la Sesión}
    Introducir la distribución Exponencial como modelo del tiempo entre eventos Poisson, demostrar la propiedad de falta de memoria, calcular el MLE del parámetro $\lambda$, y aplicar la distribución Erlang (suma de exponenciales) a sistemas de confiabilidad y teoría de colas.
  \end{block}
\end{frame}

% Slide 3: Motivación
\begin{frame}{Tiempo entre Eventos: La Exponencial como Modelo Natural}
  \footnotesize
  \begin{motivacion}[¿Por qué la distribución exponencial?]
    En el Capítulo 03 modelamos el \emph{número} de eventos en un intervalo (Poisson). Ahora modelamos el \emph{tiempo entre eventos}: si los eventos siguen un proceso de Poisson($\lambda$), los tiempos entre eventos consecutivos son i.i.d. $\text{Exp}(\lambda)$. La Exponencial es la distribución continua de mayor aplicación en ingeniería de confiabilidad, teoría de colas, y física de procesos estocásticos.
  \end{motivacion}

  \vspace{0.15cm}
  \pause
  \begin{itemize}\itemsep=0.1cm
    \item \textbf{Tiempo entre eventos:} Modela tiempos de espera hasta la primera ocurrencia de un evento Poisson. \pause
    \item \textbf{Falta de memoria:} $P(X > s + t \mid X > s) = P(X > t)$ es la propiedad definitoria. \pause
    \item \textbf{Base de la Erlang y Gamma:} Suma de $n$ exponenciales i.i.d. es Erlang$(n, \lambda)$, base de sistemas en serie.
  \end{itemize}
\end{frame}

% Slide 4: Definición Formal
\begin{frame}{Definición Formal de la Distribución Exponencial}
  \footnotesize
  Una variable aleatoria continua $X$ tiene distribución \emph{exponencial} con parámetro de tasa $\lambda > 0$ si su PDF es: \pause

  \vspace{0.1cm}
  \begin{block}{Función de Densidad Exponencial}
    $$f(x) = \begin{cases} \lambda e^{-\lambda x} & x \ge 0 \\ 0 & x < 0 \end{cases}$$
    Por normalización: $\int_{0}^{\infty} \lambda e^{-\lambda x}\,dx = \lambda \cdot \frac{1}{\lambda} = 1$. En este caso escribimos $X \sim \text{Exp}(\lambda)$.
  \end{block}

  \vspace{0.1cm}
  \pause
  \begin{alertblock}{CDF y Propiedades}
    \begin{itemize}\itemsep=0.05cm
      \item $F(x) = 1 - e^{-\lambda x}$ para $x \ge 0$; $F(0) = 0$ y $\lim_{x \to \infty} F(x) = 1$.
      \item \textbf{Cola:} $P(X > x) = e^{-\lambda x}$, forma funcional idéntica a la distribución continua de Poisson.
      \item \textbf{Media:} $\E[X] = 1/\lambda$ (interpretación: tiempo medio entre eventos).
    \end{itemize}
  \end{alertblock}
\end{frame}

% Slide 5: Esperanza y Varianza
\begin{frame}{Esperanza y Varianza: $\E[X] = 1/\lambda$, $\Var(X) = 1/\lambda^{2}$}
  \footnotesize
  Los momentos de la Exponencial tienen formas cerradas particularmente simples: \pause

  \vspace{0.1cm}
  \begin{block}{Momentos por Integración}
    \begin{align*}
      \E[X] &= \int_{0}^{\infty} x \lambda e^{-\lambda x}\,dx = \frac{1}{\lambda} \quad \text{(integración por partes)}, \\
      \E[X^{2}] &= \int_{0}^{\infty} x^{2} \lambda e^{-\lambda x}\,dx = \frac{2}{\lambda^{2}}, \\
      \Var(X) &= \E[X^{2}] - (\E[X])^{2} = \frac{2}{\lambda^{2}} - \frac{1}{\lambda^{2}} = \frac{1}{\lambda^{2}}.
    \end{align*}
  \end{block}

  \vspace{0.1cm}
  \pause
  \begin{alertblock}{MGF y Función Característica}
    \begin{itemize}\itemsep=0.05cm
      \item $M_{X}(t) = \E[e^{tX}] = \dfrac{\lambda}{\lambda - t}$ para $t < \lambda$.
      \item \textbf{Cociente de dispersión:} $\E[X] = \sigma$ (media igual a desviación estándar).
      \item \textbf{Asimetría:} $\gamma_{1} = 2$ (cola derecha pesada, hacia valores grandes).
    \end{itemize}
  \end{alertblock}
\end{frame}

% Slide 6: Propiedad de Falta de Memoria
\begin{frame}{Propiedad de Falta de Memoria: $P(X > s + t \mid X > s) = P(X > t)$}
  \footnotesize
  La Exponencial es la \emph{única} distribución continua con soporte en $[0, \infty)$ que satisface la propiedad de falta de memoria: \pause

  \vspace{0.1cm}
  \begin{block}{Demostración de la Propiedad}
    \begin{align*}
      P(X > s + t \mid X > s) = \dfrac{P(X > s + t \cap X > s)}{P(X > s)} = \dfrac{P(X > s + t)}{P(X > s)} = \dfrac{e^{-\lambda(s+t)}}{e^{-\lambda s}} = e^{-\lambda t} = P(X > t).
    \end{align*}
    El cálculo explícito de la cola exponencial cancela el factor $e^{-\lambda s}$, dejando solo $e^{-\lambda t}$.
  \end{block}

  \vspace{0.1cm}
  \pause
  \begin{alertblock}{Unicidad e Implicaciones}
    \begin{itemize}\itemsep=0.05cm
      \item \textbf{Unicidad:} Si $g(s+t) = g(s) g(t)$ con $g$ medible, entonces $g(t) = e^{-ct}$.
      \item \textbf{Aplicación:} ``El sistema no envejece'': la probabilidad de fallo futuro no depende del tiempo pasado.
      \item \textbf{Contraste:} La Uniforme NO tiene esta propiedad (probabilidad condicional depende del soporte restante).
    \end{itemize}
  \end{alertblock}
\end{frame}

% Slide 7: Distribución Erlang
\begin{frame}{Distribución Erlang: Suma de $n$ Exponenciales i.i.d.}
  \footnotesize
  La suma de $n$ variables $\text{Exp}(\lambda)$ independientes sigue la distribución \emph{Erlang}($n$, $\lambda$): \pause

  \vspace{0.1cm}
  \begin{block}{Definición y Propiedades de Erlang}
    Si $X_{1}, \dots, X_{n} \sim \text{Exp}(\lambda)$ i.i.d., entonces $S_{n} = \sum_{i=1}^{n} X_{i} \sim \text{Erlang}(n, \lambda)$ con:
    $$f_{S_{n}}(x) = \dfrac{\lambda^{n} x^{n-1} e^{-\lambda x}}{(n-1)!}, \quad x \ge 0.$$
    Momentos: $\E[S_{n}] = n/\lambda$ y $\Var(S_{n}) = n/\lambda^{2}$.
  \end{block}

  \vspace{0.1cm}
  \pause
  \begin{alertblock}{Aplicación: Sistemas en Serie}
    \begin{itemize}\itemsep=0.05cm
      \item En un sistema de $n$ componentes en serie, el tiempo hasta la falla del sistema es $T = \min(X_{1}, \dots, X_{n}) \sim \text{Exp}(n\lambda)$.
      \item En un sistema paralelo, el tiempo hasta la $k$-ésima falla es Erlang$(k, \lambda)$.
      \item Aproximación Normal: para $n \ge 30$, $S_{n} \approx N(n/\lambda, n/\lambda^{2})$ por TLC.
    \end{itemize}
  \end{alertblock}
\end{frame}

% Slide 8: Estimación por MLE
\begin{frame}{Estimación por MLE: $\hat{\lambda} = 1/\bar{X}$}
  \footnotesize
  El estimador de máxima verosimilitud para el parámetro de tasa $\lambda$ tiene una forma cerrada simple: \pause

  \vspace{0.1cm}
  \begin{block}{Derivación del MLE}
    Para $n$ muestras i.i.d. de $\text{Exp}(\lambda)$, la log-verosimilitud es:
    $$\ell(\lambda) = n \log \lambda - \lambda \sum_{i=1}^{n} x_i$$
    Derivando: $\ell'(\lambda) = n/\lambda - \sum x_i = 0$, dando $\hat{\lambda} = \dfrac{n}{\sum x_i} = \dfrac{1}{\bar{X}}$.
  \end{block}

  \vspace{0.1cm}
  \pause
  \begin{alertblock}{Propiedades Asintóticas del MLE}
    \begin{itemize}\itemsep=0.05cm
      \item \textbf{Sesgo:} $\E[\hat{\lambda}] = n\lambda/(n-1)$, así $\hat{\lambda}$ es sesgado pero consistente.
      \item \textbf{Varianza asintótica:} $\Var(\hat{\lambda}) \to \lambda^{2}/n$ vía información de Fisher $I(\lambda) = 1/\lambda^{2}$.
      \item \textbf{Normalidad asintótica:} $\sqrt{n}(\hat{\lambda} - \lambda)/\lambda \xrightarrow{d} N(0, 1)$ cuando $n \to \infty$.
    \end{itemize}
  \end{alertblock}
\end{frame}

% Slide 9: Aplicaciones
\begin{frame}{Aplicaciones en Ingeniería, Ciencias y Teoría de Colas}
  \footnotesize
  La Exponencial es la piedra angular de múltiples áreas aplicadas: \pause

  \vspace{0.15cm}
  \begin{itemize}\itemsep=0.12cm
    \item \textbf{Confiabilidad Industrial:} Tiempo hasta la falla de componentes electrónicos y mecánicos bajo la propiedad de no-envejecimiento. \pause
    \item \textbf{Teoría de Colas M/M/1:} Sistema con llegadas Poisson($\lambda$) y servicios Exp($\mu$); estable si $\rho = \lambda/\mu < 1$. \pause
    \item \textbf{Física de Procesos:} Tiempo de decaimiento radiactivo, vida media de partículas subatómicas. \pause
    \item \textbf{Ciencia de Datos:} Modelado de tiempos entre eventos en logs de sistemas, mantenimiento predictivo, y análisis de supervivencia.
  \end{itemize}
\end{frame}

% Slide 12A-1: Lab Python (Bloque 1)
\begin{frame}[fragile]{Laboratorio en Python: PDF y Falta de Memoria}
  \scriptsize
  Validación de $\int f = 1$ y demostración numérica de la propiedad de falta de memoria (\texttt{04.05\_exponential\_distribution.py}):
  \vspace{0.02cm}
  {\lstset{basicstyle=\fontsize{4.5pt}{5.5pt}\selectfont\ttfamily}
  \lstinputlisting[firstline=15, lastline=40]{../../code/04_variables_aleatorias_continuas/04.05_exponential_distribution.py}}
\end{frame}

% Slide 12A-2: Lab Python (Bloque 2)
\begin{frame}[fragile]{Laboratorio en Python: MLE y Bootstrap}
  \scriptsize
  Estimación de $\hat{\lambda} = 1/\bar{X}$, intervalos de confianza bootstrap y convergencia del MLE:
  \vspace{0.02cm}
  {\lstset{basicstyle=\fontsize{4pt}{5pt}\selectfont\ttfamily}
  \lstinputlisting[firstline=45, lastline=75]{../../code/04_variables_aleatorias_continuas/04.05_exponential_distribution.py}}
\end{frame}

% Slide 12B: Lab Python (Bloque 3)
\begin{frame}[fragile]{Laboratorio en Python: Erlang y Sistemas de Colas}
  \scriptsize
  Verificación de la distribución Erlang y aplicación a sistemas M/M/1 de teoría de colas:
  \vspace{0.02cm}
  {\lstset{basicstyle=\fontsize{4.5pt}{5.5pt}\selectfont\ttfamily}
  \lstinputlisting[firstline=85, lastline=110]{../../code/04_variables_aleatorias_continuas/04.05_exponential_distribution.py}}
\end{frame}

% Slide 12C: Salida Terminal
\begin{frame}[fragile]{Laboratorio en Python: Salida en Terminal}
  \scriptsize
  Salida estándar generada al ejecutar el script del laboratorio en Python:
  \vspace{0.02cm}
  \begin{block}{Salida Estándar en Consola}
    \fontsize{4pt}{5pt}\selectfont
    \begin{verbatim}
=== Block 1: PDF & Memoryless Property ===
Exp(2.0) PDF integral: 1.00000000
Memoryless: P(X>s+t|X>s) = P(X=t) (diff < 1e-15)
Moments: E[X]=0.5 | E[X^2]=0.5 | Var=0.25

=== Block 2: MLE & Quantile Estimation ===
True lambda=2.0, MLE=2.3642, Bootstrap SE=0.3699
95% CI: [1.8213, 3.2714] (contains true value)
Convergence: n=5000: |error|=0.0052

=== Block 3: Erlang & Reliability ===
Sum of 5 Exp(1) = Erlang(5,1): mean=5.006 | KS D=0.003
Reliability 10 components: MTTF=100h, P(T>100)=0.367
M/M/1 queue: rho=0.833, P(busy)=83.3%, L=5.0
    \end{verbatim}
  \end{block}
\end{frame}

% Slide 13A: Ejercicio Nivel 1 (Enunciado)
\begin{frame}{Ejercicio en Clase --- Nivel 1: Fundamental (1/2)}
  \footnotesize
  \begin{block}{Problema 4.5.1 --- PDF y CDF Exponencial (Enunciado)}
    Sea $X \sim \text{Exp}(\lambda = 2)$ con PDF $f(x) = 2e^{-2x}$ para $x \ge 0$. Verifique que $\int_{0}^{\infty} f(x)\,dx = 1$, calcule la CDF $F(x) = 1 - e^{-2x}$, y determine $P(X \le 1)$ y $P(X \ge 3)$.
  \end{block}

  \vspace{0.15cm}
  \pause
  \begin{alertblock}{Planteamiento del Modelo}
    \begin{itemize}\itemsep=0.04cm
      \item Soporte: $[0, \infty)$ (semi-acotado, cola derecha).
      \item PDF exponencial: $f(x) = \lambda e^{-\lambda x}$ con tasa $\lambda = 2$.
      \item CDF: $F(x) = 1 - e^{-\lambda x}$ y cola: $P(X > x) = e^{-\lambda x}$.
    \end{itemize}
  \end{alertblock}
\end{frame}

% Slide 13B: Ejercicio Nivel 1 (Resolución)
\begin{frame}{Ejercicio en Clase --- Nivel 1: Fundamental (2/2)}
  \scriptsize
  Verificamos la normalización y calculamos las probabilidades: \pause

  \vspace{0.05cm}
  \begin{block}{Normalización y CDF}
    \begin{align*}
      \int_{0}^{\infty} 2e^{-2x}\,dx = 2 \cdot \left[ -\dfrac{e^{-2x}}{2} \right]_{0}^{\infty} = 2 \cdot (0 - (-\tfrac{1}{2})) = 1.
    \end{align*}
    La CDF es $F(x) = 1 - e^{-2x}$ para $x \ge 0$.
  \end{block}

  \vspace{0.05cm}
  \pause
  \begin{alertblock}{Probabilidades}
    \scriptsize
    \begin{itemize}\itemsep=0.04cm
      \item $P(X \le 1) = 1 - e^{-2} \approx 1 - 0.1353 = 0.8647$.
      \item $P(X \ge 3) = e^{-6} \approx 0.00248$.
    \end{itemize}
    Interpretación: el 86.5\% de las observaciones caen en $[0, 1]$, pero solo el 0.25\% exceden 3. La cola exponencial decae rápidamente.
  \end{alertblock}
\end{frame}

% Slide 14A: Ejercicio Nivel 2 (Enunciado)
\begin{frame}{Ejercicio en Clase --- Nivel 2: Operativo (1/2)}
  \footnotesize
  \begin{block}{Problema 4.5.5 --- Cuantiles de la Exponencial (Enunciado)}
    Sea $X \sim \text{Exp}(\lambda = 0.5)$ (media 2 minutos). Calcule los cuantiles $q_{0.10}$, $q_{0.50}$ (mediana) y $q_{0.90}$, y encuentre el cuantil $q_{0.95}$ en el contexto de un sistema de colas.
  \end{block}

  \vspace{0.15cm}
  \pause
  \begin{alertblock}{Planteamiento}
    \scriptsize
    La CDF es $F(x) = 1 - e^{-0.5 x}$, así $q_{p} = -\frac{1}{0.5}\ln(1-p) = -2\ln(1-p)$. El cuantil $q_{0.95}$ define el tiempo de espera bajo el cual ocurre el 95\% de los clientes.
  \end{alertblock}
\end{frame}

% Slide 14B: Ejercicio Nivel 2 (Resolución)
\begin{frame}{Ejercicio en Clase --- Nivel 2: Operativo (2/2)}
  \scriptsize
  Aplicamos la fórmula de cuantiles $q_{p} = -2\ln(1-p)$: \pause

  \vspace{0.05cm}
  \begin{block}{Cálculo de Cuantiles}
    \begin{itemize}\itemsep=0.05cm
      \item $q_{0.10} = -2\ln(0.9) \approx 0.211$ min. \pause
      \item $q_{0.50} = 2\ln 2 \approx 1.386$ min (mediana). \pause
      \item $q_{0.90} = -2\ln(0.1) \approx 4.605$ min. \pause
      \item $q_{0.95} = -2\ln(0.05) \approx 5.991$ min: el 95\% de los clientes espera menos de 6 minutos.
    \end{itemize}
  \end{block}

  \vspace{0.05cm}
  \pause
  \begin{alertblock}{Decisión de Gestión}
    \scriptsize
    Para garantizar que el 95\% de los clientes espere menos de 6 minutos, el sistema debe tener capacidad de servicio $\mu > 1/6 \approx 0.167$ clientes/seg. Si $\lambda = 0.5$ clientes/seg, se necesitan al menos 3 servidores para estabilidad operativa.
  \end{alertblock}
\end{frame}

% Slide 15A: Ejercicio Nivel 3 (Enunciado)
\begin{frame}{Ejercicio en Clase --- Nivel 3: Analítico (1/2)}
  \footnotesize
  \begin{block}{Problema 4.5.7 --- Deducción del MLE Exponencial (Enunciado)}
    Sea $X_{1}, \dots, X_{n}$ i.i.d. de $\text{Exp}(\lambda)$. Deduzca el MLE del parámetro $\lambda$, demuestre que $\hat{\lambda} = 1/\bar{X}$, y calcule el sesgo y la varianza asintótica.
  \end{block}

  \vspace{0.1cm}
  \pause
  \begin{alertblock}{Estrategia de Deducción}
    \scriptsize
    Escriba la log-verosimilitud $\ell(\lambda) = n \log \lambda - \lambda \sum x_i$, derive respecto a $\lambda$, iguale a cero. Para varianza asintótica, use la información de Fisher $I(\lambda) = 1/\lambda^{2}$.
  \end{alertblock}
\end{frame}

% Slide 15B: Ejercicio Nivel 3 (Resolución)
\begin{frame}{Ejercicio en Clase --- Nivel 3: Analítico (2/2)}
  \scriptsize
  Derivamos la log-verosimilitud: \pause

  \vspace{0.05cm}
  \begin{block}{Demostración}
    \scriptsize
    \begin{align*}
      \ell(\lambda) &= \sum_{i=1}^{n} \log f(x_i \mid \lambda) = n \log \lambda - \lambda \sum_{i=1}^{n} x_i. \\
      \ell'(\lambda) &= \dfrac{n}{\lambda} - \sum_{i=1}^{n} x_i = 0 \implies \hat{\lambda} = \dfrac{n}{\sum x_i} = \dfrac{1}{\bar{X}}.
    \end{align*}
    Sesgo: $\E[\hat{\lambda}] = n\lambda/(n-1) > \lambda$. Consistencia: $\hat{\lambda} \xrightarrow{p} \lambda$. La varianza asintótica es $\Var(\hat{\lambda}) \to \lambda^{2}/n$ por la información de Fisher. \quad \qedhere
  \end{block}
\end{frame}

% Slide 16A: Ejercicio Nivel 4 (Enunciado)
\begin{frame}{Ejercicio en Clase --- Nivel 4: Desafiante (1/2)}
  \footnotesize
  \begin{block}{Problema 4.5.9 --- Sistema de Confiabilidad en Serie (Enunciado)}
    En un sistema de confiabilidad con $n$ componentes en serie, si los tiempos hasta la falla son i.i.d. $X_{i} \sim \text{Exp}(\lambda)$ independientes, demuestre que el tiempo de falla del sistema $T = \min(X_{1}, \dots, X_{n})$ sigue $\text{Exp}(n\lambda)$.
  \end{block}

  \vspace{0.15cm}
  \pause
  \begin{alertblock}{Aplicación}
    \scriptsize
    Para $n = 10$ componentes con $\lambda = 0.001$ fallas/hora, calcule el tiempo medio hasta la falla del sistema, $P(T > 100)$ y la mediana del tiempo de falla.
  \end{alertblock}
\end{frame}

% Slide 16B: Ejercicio Nivel 4 (Resolución)
\begin{frame}{Ejercicio en Clase --- Nivel 4: Desafiante (2/2)}
  \scriptsize
  Demostramos vía independencia y la propiedad de falta de memoria: \pause

  \vspace{0.05cm}
  \begin{block}{Demostración Formal}
    \scriptsize
    Por independencia de los $X_{i}$:
    \begin{align*}
      P(T > t) = P(\min(X_{1}, \dots, X_{n}) > t) = \prod_{i=1}^{n} P(X_{i} > t) = \prod_{i=1}^{n} e^{-\lambda t} = e^{-n\lambda t}.
    \end{align*}
    Esto es exactamente la cola de $\text{Exp}(n\lambda)$, así $T \sim \text{Exp}(n\lambda)$. \quad \qedhere
  \end{block}

  \vspace{0.05cm}
  \pause
  \begin{alertblock}{Aplicación Numérica}
    \scriptsize
    Con $n = 10$ y $\lambda = 0.001$: $\E[T] = 1/(n\lambda) = 100$ horas. $P(T > 100) = e^{-1} \approx 0.368$. La mediana es $m = \ln 2/(n\lambda) = 69.3$ horas $\approx 2.89$ días. El sistema tiene un 36.8\% de probabilidad de sobrevivir más de 100 horas.
  \end{alertblock}
\end{frame}

% Slide 17: Cierre
\begin{frame}{Cierre de Sección y Síntesis sobre la Distribución Exponencial}
  \begin{reflexion}[La Exponencial como Distribución sin Memoria]
    La distribución Exponencial es la única distribución continua con soporte en $[0, \infty)$ que satisface la propiedad de falta de memoria. Esta propiedad, derivada del proceso de Poisson subyacente, la hace única en confiabilidad, teoría de colas, y análisis de supervivencia. Su MLE tiene forma cerrada ($\hat{\lambda} = 1/\bar{X}$), y la suma de $n$ exponenciales i.i.d. produce la distribución Erlang.
  \end{reflexion}

  \vspace{0.2cm}
  \pause
  \begin{block}{Lecciones Clave de la Distribución Exponencial}
    \begin{itemize}\itemsep=0.08cm
      \item \textbf{PDF:} $f(x) = \lambda e^{-\lambda x}$ modela tiempos de espera continuos.
      \item \textbf{Falta de memoria:} $P(X > s + t \mid X > s) = P(X > t)$ es la propiedad definitoria (y única).
      \item \textbf{Erlang:} La suma de $n$ exponenciales i.i.d. es Erlang$(n, \lambda)$, base de sistemas en serie.
    \end{itemize}
  \end{block}
\end{frame}

% Slide 18: Perspectiva Modular
\begin{frame}{Perspectiva Modular --- El Próximo Reto}
  \scriptsize
  \begin{retoclase}[Para la próxima sesión: Distribución Normal y Puntaje $Z$]
    La Exponencial modela tiempos de espera. Ahora estudiaremos la distribución más importante de toda la estadística: la \emph{Normal} $N(\mu, \sigma^{2})$, cuya PDF en forma de campana modela errores de medición, ruido gaussiano, y muchos fenómenos naturales. La estandarización al puntaje $Z = (X - \mu)/\sigma$ permite calcular probabilidades usando tablas universales.
  \end{retoclase}

  \vspace{0.08cm}
  \pause
  \begin{alertblock}{Preparación Recomendada}
    \begin{itemize}\itemsep=0.06cm
      \item Repasa el Teorema del Límite Central y la convergencia de sumas de variables i.i.d. a la Normal.
      \item Reflexiona sobre por qué la Normal aparece en tantos fenómenos naturales a pesar de la rareza de sus colas.
    \end{itemize}
  \end{alertblock}
\end{frame}

\end{document}
